\section*{Hoofdstuk \ref{ch:g11:organic-skeletons} --- Organische moleculen: skeletten en namen}

\begin{solution}{exo:g11:organic-skeletons:1}
Beide hebben zes koolstofatomen (uiteinden en hoekpunten) en
\ce{C6H14}: hexaan en 2-methylpentaan zijn \omterm{def:g11:organic-skeletons:isomer}{isomeren}.
\end{solution}

\begin{solution}{exo:g11:organic-skeletons:2}
Propaan; hexaan; 2-methylbutaan.
\end{solution}

\begin{solution}{exo:g11:organic-skeletons:3}
\ce{CH3-CH2-CH3}; \chemfig{CH_3-CH(-[2]CH_3)-CH_3};
\chemfig{H-C(-[2]H)(-[6]H)-C(-[2]H)(-[6]H)-H}.
\end{solution}

\begin{solution}{exo:g11:organic-skeletons:4}
Het eerste, 3-methylpentaan: 6 koolstofatomen; de \ce{CH3}-uiteinden
(drie ervan) dragen 9 waterstofatomen, de twee \ce{CH2}-hoekpunten 4 en
het \ce{CH}-vertakkingspunt 1: \ce{C6H14}. Cyclohexaan: 6 hoekpunten, elk
\ce{CH2}: \ce{C6H12}.
\end{solution}

\begin{solution}{exo:g11:organic-skeletons:5}
\ce{C8H18}. $2n + 2 = 20$ geeft $n = 9$: nonaan, \ce{C9H20}.
\end{solution}

\begin{solution}{exo:g11:organic-skeletons:6}
\setchemfig{atom sep=1.6em}
3-methylhexaan \chemfig{-[:30]-[:-30]-[:30](-[:90])-[:-30]-[:30]-[:-30]},
\ce{C7H16}; 2,2-dimethylbutaan \chemfig{-[:0](-[:90])(-[:-90])-[:0]-[:30]},
\ce{C6H14}; 3-ethylpentaan
\chemfig{-[:30]-[:-30](-[:-90]-[:-30])-[:30]-[:-30]}, \ce{C7H16}.
\end{solution}

\begin{solution}{exo:g11:organic-skeletons:7}
Pentaan, 2-methylbutaan en 2,2-dimethylpropaan, getekend in de figuur
van het hoofdstuk.
\end{solution}

\begin{solution}{exo:g11:organic-skeletons:8}
``2-ethylbutaan'' heeft een langste keten van vijf koolstofatomen door de
ethylgroep: 3-methylpentaan. ``4-methylpentaan'' is genummerd vanaf het
verkeerde uiteinde: 2-methylpentaan. ``1-methylbutaan'' is gewoon een
keten van vijf koolstofatomen: pentaan.
\end{solution}

\begin{solution}{exo:g11:organic-skeletons:9}
2,2-dimethylpropaan (\qty{9.5}{\celsius}) < 2-methylbutaan
(\qty{27.8}{\celsius}) < pentaan (\qty{36.1}{\celsius}). Dezelfde \omterm{def:g7:atoms-and-molecules:atom}{atomen}
zijn steeds compacter geschikt: minder contact tussen \omterm{def:g7:atoms-and-molecules:molecule}{moleculen},
zwakkere \omterm{def:g11:polarity-and-cohesion:van-der-waals}{vanderwaalsbindingen}, lager kookpunt.
\end{solution}

\begin{solution}{exo:g11:organic-skeletons:10}
Als de keten zich tot een ring sluit, ontstaat er één \ce{C-C}-binding
meer, die de plaats inneemt van twee \ce{C-H}-bindingen (één aan elk
uiteinde). Verschillende \omterm{def:g11:organic-skeletons:formulas}{molecuulformules}: geen \omterm{def:g11:organic-skeletons:isomer}{isomeren}. Cyclohexaan is
geen \omterm{def:g11:organic-skeletons:alkane}{alkaan} (het skelet is een ring; de formule is niet
\ce{C_{n}H_{2n+2}}).
\end{solution}

\begin{solution}{exo:g11:organic-skeletons:11}
Butaan en 2-methylpropaan (\ce{C4H10}). Propaan is \ce{C3H8} en
cyclobutaan \ce{C4H8}.
\end{solution}

\begin{solution}{exo:g11:organic-skeletons:12}
\setchemfig{atom sep=1.5em}
\begin{center}
\small
\begin{tabular}{ccc}
\chemfig{-[:30]-[:-30]-[:30]-[:-30]-[:30]} &
\chemfig{-[:30](-[:90])-[:-30]-[:30]-[:-30]} &
\chemfig{-[:30]-[:-30](-[:-90])-[:30]-[:-30]} \\
hexaan & 2-methylpentaan & 3-methylpentaan \\[6pt]
\chemfig{-[:0](-[:90])(-[:-90])-[:0]-[:30]} &
\chemfig{-[:30](-[:90])-[:-30](-[:-90])-[:30]} & \\
2,2-dimethylbutaan & 2,3-dimethylbutaan &
\end{tabular}
\end{center}
\end{solution}

\begin{solution}{exo:g11:organic-skeletons:13}
Twee lange zigzags kunnen over hun hele lengte naast elkaar liggen, met
veel \omterm{def:g7:atoms-and-molecules:atom}{atomen} die elkaar raken; twee ballen raken elkaar maar op een klein
stukje. \omterm{def:g11:polarity-and-cohesion:van-der-waals}{Vanderwaalsbindingen} werken tussen \omterm{def:g7:atoms-and-molecules:atom}{atomen} die elkaar raken: ze
zijn groter tussen onvertakte \omterm{def:g7:atoms-and-molecules:molecule}{moleculen}, die een hogere temperatuur nodig
hebben om van elkaar los te komen.
\end{solution}

\begin{solution}{exo:g11:organic-skeletons:14}
De langste keten heeft zes koolstofatomen; vanaf links genummerd staan de
substituenten op 2, 4, 4, vanaf rechts op 3, 3, 5; het eerste verschil
(2 tegen 3) beslist: 2,4,4-trimethylhexaan.
\end{solution}

\begin{solution}{exo:g11:organic-skeletons:15}
\setchemfig{atom sep=1.6em}
\chemfig{-[:0](-[:90])(-[:-90])-[:-30]-[:30](-[:90])-[:-30]}: een keten van
vijf koolstofatomen met twee methylgroepen op koolstofatoom 2 en één op
koolstofatoom 4, in totaal acht koolstofatomen; $2 \times 8 + 2 = 18$
waterstofatomen, \ce{C8H18}, de formule van octaan: het zijn \omterm{def:g11:organic-skeletons:isomer}{isomeren}.
\omterm{def:g11:organic-skeletons:formulas}{Beknopte structuurformule}:
\chemfig{CH_3-C(-[2]CH_3)(-[6]CH_3)-CH_2-CH(-[2]CH_3)-CH_3}.
\end{solution}

\begin{solution}{pb:g11:organic-skeletons:1}
\setchemfig{atom sep=1.5em}
\textbf{1.} \ce{C4H10}: met $n = 4$ is $2n + 2 = 10$.

\textbf{2.} \chemfig{H-C(-[2]H)(-[6]H)-C(-[2]H)(-[6]H)-C(-[2]H)(-[6]H)-C(-[2]H)(-[6]H)-H}.

\textbf{3.} \ce{CH3-CH2-CH2-CH3}; \chemfig{-[:30]-[:-30]-[:30]}.

\textbf{4.} \chemfig{H-C(-[2]H)(-[6]H)-C(-[6]H)(-[2,1.5]C(-[1]H)(-[3]H)(-[2]H))-C(-[2]H)(-[6]H)-H};
\chemfig{CH_3-CH(-[2]CH_3)-CH_3}; \chemfig{-[:30](-[:90])-[:-30]}. Het
heeft dezelfde \omterm{def:g7:atoms-and-molecules:atom}{atomen}, vier koolstofatomen en tien waterstofatomen, in een
andere volgorde aan elkaar gebonden.

\textbf{5.} Ja: \omterm{def:g11:organic-skeletons:isomer}{structuurisomeren} (verschillende skeletten).

\textbf{6.} Alle drie: hun kookpunten liggen onder \qty{20}{\celsius}.

\textbf{7.} Bij \qty{-5}{\celsius}, onder zijn kookpunt, blijft butaan
bij normale druk vloeibaar en geeft het weinig gas af.

\textbf{8.} Propaan (\qty{-42}{\celsius}) en 2-methylpropaan
(\qty{-11.7}{\celsius}) koken onder wintertemperaturen: ze geven in de
kou nog genoeg gas.

\textbf{9.} 2-methylpropaan is vertakt en compacter: zwakkere
\omterm{def:g11:polarity-and-cohesion:van-der-waals}{vanderwaalsbindingen}.

\textbf{10.} Met drie koolstofatomen is het enig mogelijke skelet de
keten \ce{C-C-C}: elk koolstofatoom dat je ergens anders zet, geeft
dezelfde keten.

\textbf{11.} 2-methylbutaan.

\textbf{12.} De keten moet genummerd worden vanaf het uiteinde dat het
dichtst bij de substituent ligt: 2, niet 3.

\textbf{13.} 2, 3 en 5.

\textbf{14.} Heptaan, \chemfig{-[:30]-[:-30]-[:30]-[:-30]-[:30]-[:-30]}.

\textbf{15.} 2-methylhexaan \chemfig{-[:30](-[:90])-[:-30]-[:30]-[:-30]-[:30]}
en 3-methylhexaan \chemfig{-[:30]-[:-30](-[:-90])-[:30]-[:-30]-[:30]}.
``4-methylhexaan'' is 3-methylhexaan, genummerd vanaf het verkeerde
uiteinde.

\textbf{16.} 2,2-dimethylpentaan, 2,3-dimethylpentaan,
2,4-dimethylpentaan en 3,3-dimethylpentaan.

\textbf{17.} Ja, 3-ethylpentaan. Op koolstofatoom 2 zou de ethylgroep
een keten van zes koolstofatomen door die groep vormen: het \omterm{def:g7:atoms-and-molecules:molecule}{molecuul} zou
3-methylhexaan zijn.

\textbf{18.} 2,2,3-trimethylbutaan,
\chemfig{-[:0](-[:90])(-[:-90])-[:0](-[:90])-[:0]}.

\textbf{19.} $1 + 2 + 4 + 1 + 1 = 9$.

\textbf{20.} Heptaan heeft 9 \omterm{def:g11:organic-skeletons:isomer}{structuurisomeren}, heptaan zelf
meegerekend.
\end{solution}
